\section{Introduction}
\label{sec:intro}

A familiar shape recurs across the history of delegated emergency authority.
A legislature, faced with a circumstance it characterizes as exceptional,
grants an executive officer a power that was not previously available within the ordinary structure of the legal order.
The grant is framed as time-bounded.
The authority is to be used only while the exceptional circumstance persists,
and is to lapse, by its own terms, once the circumstance abates.
The grant is sometimes accompanied by a sunset provision,
sometimes by a periodic reauthorization requirement,
sometimes by reporting obligations to a legislative committee,
and sometimes by judicial review of the executive's exercise of the granted authority.
In each case, the structural intuition motivating the grant is that the authority is borrowed,
not given outright;
that it must return to the legislature when the exception ends;
and that the legal order will resume its ordinary shape thereafter.

The structural intuition is, with regularity that does not appear to be coincidental, falsified.
Article~48 of the Weimar Constitution authorized the Reich President to take ``the necessary measures''
to restore public security in cases where it was ``substantially disturbed or endangered''
\cite{ref_weimar_const_art48}.
Article~48 required notification to the Reichstag and permitted the Reichstag to demand the suspension of any measure taken under it.
In its first twelve years of operation, the authority was invoked approximately 250 times \cite{ref_kershaw_hitler}.
By the time of the Reichstag Fire Decree in February 1933, the authority had become a routine instrument of governance
rather than the exceptional remedy the constitutional text described \cite{ref_caldwell_popular_sovereignty}.
The transition from emergency power to instrument of dictatorship occurred through a sequence of contested political decisions,
each conducted within a constitutional grammar that did not require lapse-by-construction.
There was no single moment at which the authority's character changed.
There was, instead, a sequence of admissions, each of which appeared, at the time of its admission,
to be a defensible exercise of a power constitutionally available to its holder.

The Authorization for Use of Military Force enacted on September 18, 2001 \cite{ref_aumf_2001}
permits the President to use ``all necessary and appropriate force'' against those determined to have
authored or aided the September 11 attacks.
The text contains no sunset.
It contains no geographic limit.
It is, twenty-five years later, the primary legal authority for armed action by United States forces
in countries that did not exist as targets of policy when the authorization was enacted
\cite{ref_brennan_center_aumf,ref_chesney_aumf}.
The authorization has been cited as legal basis in at least eight named theaters by 2024
\cite{ref_jaffer_aumf}.
Periodic legislative efforts to repeal and replace the 2001 AUMF have not produced enacted legislation.
The pattern admits both a political and a structural reading.
A run of bipartisan repeal-and-replace proposals,
including Kaine-Young, Lee-Murphy, and the 2023 repeal of the 1991 and 2002 Iraq AUMFs that left the 2001 AUMF intact,
demonstrates that the political coalition for repeal has at times assembled
\cite{ref_kaine_young_aumf,ref_iraq_aumf_repeal_2023}.
What has not converged is agreement on a successor authority.
The structural reading this Article advances does not displace the political account;
it identifies an additional feature of the original grant that lowers the cost of continuation
relative to replacement at each successive period.

Section~702 of FISA, codified at 50 U.S.C. \S~1881a, governs programmatic acquisition of foreign-intelligence information
from non-US persons reasonably believed to be located abroad.
The program is not itself an emergency authority,
but \S~1881a(c)(2) contains a sub-provision permitting the Attorney General and Director of National Intelligence
to authorize targeting prior to FISC approval of a new certification.
That sub-provision, and its interaction with the broader \S~702 retention and querying framework documented in
declassified FISC opinions and PCLOB reports
\cite{ref_pclob_702_report,ref_fisc_702_opinions},
exhibits the structural pattern this Article identifies.
The structural defect is a grant of authority intended for time-bounded use that contains no
mechanism by which the underlying surveillance is rolled back when the emergency lapses.
The information collected during the emergency window remains in the holdings of the collecting agency.
The targets do not return to a prior state.

Article~17 of the General Data Protection Regulation grants data subjects a right to erasure of their personal data
in specified circumstances \cite{ref_gdpr_art17}.
The Court of Justice of the European Union's enforcement of Article~17 has produced erasure orders against
search-engine providers and platform operators that, in form, are time-bounded executive acts of administrative agencies.
In practice, an erasure order does not contemplate a path back to the prior state.
A search-engine index from which a result has been removed is not, on the regulation's default operation, restored,
although the underlying balance is re-evaluable and the controller may re-index where the supporting facts change
\cite{ref_kuner_gdpr_erasure}.
The order is structurally irreversible at the moment of its admission,
even though no statutory text characterizes it as such.

Section~230 of the Communications Decency Act is, in its primary form, a liability shield rather than a takedown
statute: \S~230(c)(1) immunizes interactive computer services from being treated as the publisher of third-party content,
and \S~230(c)(2) protects good-faith restriction of access to objectionable material
\cite{ref_section230_shield}.
The formal notice-and-takedown regime in United States law is the Digital Millennium Copyright Act's \S~512 framework,
which includes a \S~512(g) counter-notice procedure that obligates the service provider to restore content
within a statutorily specified window absent the filing of suit by the original complainant
\cite{ref_dmca_512,ref_dmca_512g_counter_notice}.
Section 230's incentive shadow produces takedown-like behavior in platform self-regulation,
though the statute itself does not authorize takedown.
The structural shape of the time-bounded executive act,
on this Article's reading, locates within the DMCA \S~512 regime and within the broader pattern of platform self-regulation
that Section 230 shelters \cite{ref_keller_takedowns}.
That characterization is one this Article advances against the consensus of the Section 230 academy
(Goldman, Citron, Keller),
which has, in differing directions, resisted treating Section 230 as a takedown regime
\cite{ref_goldman_section230_2019,ref_keller_overremoval_2015}.
The removed content, once removed, is not statutorily required to be restored upon any later event,
and the \S~512(g) counter-notice path, although a partial rollback mechanism, is in empirical practice underutilized.

These five regimes do not share doctrinal substance.
The constitutional emergency power of a Weimar-era Reich President bears no doctrinal relation to the
content-moderation practices of a contemporary American social-media platform,
and the surveillance authorities of FISA Section~702 bear no doctrinal relation to either.
They share, however, a structural feature that this Article identifies and names:
the original grant of authority was made without a typed rollback witness.
The authority could be admitted without exhibiting, at the moment of admission,
a constructible path back to the state of affairs that existed before the authority was exercised.

The Article's central claim is that this missing structural element is not a peripheral defect.
It is the explanation for the ratcheting failure mode that has, in each of these regimes and others,
transformed a time-bounded emergency authority into a permanent feature of the legal order.
The claim is not that a typed rollback witness, had it been required at the time these authorities were enacted,
would have prevented every adverse exercise of the authorities.
It is the narrower and more defensible claim that the absence of the typed rollback witness made the ratcheting
pattern structurally available, and that its presence would have made the same pattern structurally unavailable,
although the political pressure that produced the underlying decrees would have remained.

The grammar that supplies the typed rollback witness is the grammar of bounded executive action,
developed in a companion technical paper that constructs a Lean-attestable substrate for receipt-bounded
polities \cite{ref_chio_parent_paper}.
The companion paper's contribution is the formal substrate.
This Article's contribution is the application of the substrate's grammar to the doctrinal pattern of
delegated emergency authority,
the demonstration that the absence of the grammar is the structural explanation for the ratcheting pattern,
and the exhibition of one implementation drawn from the substrate.

Part~\ref{sec:pattern} develops the historical pattern in detail and situates it within the Schmitt-Agamben
tradition on the state of exception.
Part~\ref{sec:cases} presents five case studies, with explicit acknowledgement that the case studies vary in
how cleanly they fit the framing.
Part~\ref{sec:grammar} states the formal grammar and explains the role of the typed rollback witness.
Part~\ref{sec:substrate} cites the Lean substrate without re-deriving it.
Part~\ref{sec:implementation} exhibits one implementation as a working example.
Part~\ref{sec:limits} acknowledges the limits of the formal approach,
particularly the irreducibly political content of the exception that the wrapper cannot capture.
Part~\ref{sec:conclusion} concludes.
